Il Linguaggio di Programmazione Java permette di semplificare ed ottimizzare l'esecuzione parallela di differenti \textit{Task} mettendo a disposizione del Programmatore il Costrutto delle \textit{Thread Pool}, illustrato nel presente Capitolo.

\section{\textit{Thread Pools}}
Si tratta di gruppi di \textit{Thread} cosiddetti ``\textit{Worker}'' che attendono l'allocazione di \textit{Task} da eseguire\cite{ThreadPool_JavaTutorial}. \par
Questo approccio è spesso utilizzato in quanto l'allocazione di un \textit{Thread} richiede un \textit{Overhead} di risorse che potrebbe divenire non banale con il crescere del numero di questi. Altro vantaggio portato da una filosofia che privilegi il riutilizzo di \textit{Thread} è quella di poter limitare il tempo utilizzato per la loro creazione e distruzione. \par
Non essendo però più necessariamente in grado di assegnare un singolo \textit{Thread} a ogni \textit{Task} dell'utente, il grado di parallelismo raggiungibile implementando una soluzione come questa va ovviamente a ridursi. \par
Ogni \textit{Thread Pool} dispone di una lista di \textit{Task} da eseguire da cui ciascun \textit{Thread} andrà a scegliere. Una volta che un \textit{Thread} termina l'esecuzione di un \textit{Task}, esso andrà a sceglierne un altro dalla lista fin tanto che ve ne saranno presenti. A seguire è illustrato un esempio.
\begin{figure}[H]
    \begin{center}
        \includesvg[width=\textwidth]{images/assets/ThreadPool1.svg}
        \label{fig:ThreadPool1}
        \caption{Allocazione iniziale dei \textit{Task} per ogni \textit{Thread} della \textit{Pool}. Immagine creata con l'ausilio dell'editor Draw.io\cite{drawio}.}
    \end{center}
\end{figure}
\begin{figure}[H]
    \begin{center}
        \includesvg[width=\textwidth]{images/assets/ThreadPool2.svg}
        \label{fig:ThreadPool2}
        \caption{Possibile allocazione dei \textit{Task} dopo la terminazione dell'esecuzione del \textit{Task} numero 3 e del \textit{Task} numero 1. Immagine creata con l'ausilio dell'editor Draw.io\cite{drawio}.}
    \end{center}
\end{figure}

\section{Implementazione di \textit{Thread Pools}}
L'implementazione di una \textit{Thread Pool} può avvenire in maniere differenti a seconda delle necessità.
\subsection{La Classe \texttt{Executors}}
Il Metodo \texttt{newFixedThreadPool(int numOfThreads)} della Classe \texttt{Executors}\cite{Executors_JavaDocs} permette di creare una \textit{Thread Pool} con un numero fissato di \textit{Worker Thread} atti all'esecuzione dei \textit{Task} loro assegnati. In ogni momento al massimo \texttt{numOfThreads} \textit{Thread} saranno attivi e in esecuzione sulla \textit{Pool}. In caso di errori durante l'esecuzione di un \textit{Task}, il \textit{Thread} assegnatogli ne sceglierà un altro.
\subsection{La Classe \texttt{ThreadPoolExecutor}}
La Classe \texttt{ThreadPoolExecutor}\cite{ThreadPoolExecutor_JavaDocs} consente di creare \textit{Thread Pools} più complesse, in quanto dispone di maggiori parametri per adattarsi all'utilizzo in condizioni differenti: è possibile infatti scegliere il numero minimo e massimo di \textit{Thread} che possono essere allocati in ogni momento a seconda del carico. \par
La scelta di un corretto numero di \textit{Thread} diviene ora di fondamentale importanza al fine di non sprecare risorse o di allocarne troppo poche. \par
È possibile utilizzarla in congiunzione con l'Interfaccia \texttt{ExecutorService}\cite{ExecutorService_JavaDocs} (che estende) per ottenere oggetti di tipo \texttt{Future}\cite{Future_JavaDocs} e controllare il progresso dei \textit{Task} sottomessi.
\subsection{La Classe \texttt{ScheduledThreadPoolExecutor}}
Quest'ultima è un'estensione della Classe precedente con la differenza che mentre prima i \textit{Task} venivano eseguiti non appena possibile, ora è possibile schedularne l'inizio dell'esecuzione\cite{ScheduledThreadPoolExecutor_JavaDocs}. Va utilizzata con l'Interfaccia \texttt{ScheduledExecutorService}\cite{ScheduledExecutorService_JavaDocs}.
\subsection{La Classe \texttt{ForkJoinPool}}
È un'alternativa radicalmente diversa dalle altre presentate in quanto si basa sulla tecnica del \textit{Work-Stealing}\cite{ForkJoinPool_JavaDocs}, ovvero i \textit{Worker Thread} cercheranno di eseguire \textit{Task} sottomesse alla \textit{Pool} dal Programma o da altri \textit{Task} come risultato di una scissione di una scissione in unità più piccole. Anch'essa estende l'Interfaccia \texttt{ExecutorService}.
\subsection{Primitive di \texttt{ExecutorService}}
L'Interfaccia \texttt{ExecutorService} ricopre un ruolo centrale nell'implementazione di \textit{Thread Pool} dispone di Metodi per la gestione dell'esecuzione dei vari \textit{Task} assegnati ad esse. Se ne riportano i più importanti:
\begin{itemize}
    \item \textbf{Il metodo \texttt{execute()}} \par
    Metodo ereditato dall'Interfaccia \texttt{Executor}\cite{Executor_JavaDocs}, permette di avviare l’esecuzione di un Oggetto di tipo \texttt{Runnable} non appena possibile.
    \item \textbf{Il metodo \texttt{submit()}} \par
    Permette di inserire nella Coda un nuovo \textit{Task}: esso può essere sia di tipo \texttt{Runnable} - quindi senza Valore di Ritorno - che \texttt{Callable}, e quindi con un Valore di Ritorno.
    \item \textbf{Il metodo \texttt{shutdown()}} \par
    Attende la terminazione di tutti i \textit{Task} e termina i \textit{Thread} della \textit{Pool} per poi distruggerla, non accettando nuovi inserimenti.
    \item \textbf{Il metodo \texttt{shutdownNow()}} \par
    Tenta di terminare i \textit{Task} in esecuzione, cancella l'esecuzione dei prossimi in coda e non accetta nuovi inserimenti su di essa, ritornando una lista di \textit{Task} che attendevano di eseguire al momento dell'invocazione.
    \item \textbf{Il metodo \texttt{invokeAny()}} \par
    Esegue i \textit{Task} indicati dalla lista passata come parametro al Metodo stesso, ritornando il risultato del primo \textit{Thread} che esegue correttamente il \textit{Task}.
    \item \textbf{Il metodo \texttt{invokeAll()}} \par
    Esegue i \textit{Task} indicati dalla lista passata come parametro al Metodo stesso, ritornando una lista di Oggetti di tipo \texttt{Future} che contengano lo stato e il risultato dell'esecuzione una volta che tutti hanno terminato.
    \item \textbf{Il metodo \texttt{awaitTermination()}} \par
    Blocca l'esecuzione del \textit{Thread} principale fin tanto che l'esecuzione dei \textit{Task} non termina.
\end{itemize}
\subsection{Esempio di Implementazione di una \textit{Thread Pool} tramite \texttt{ExecutorService}}
Si riporta a seguire un esempio di implementazione di una \textit{ThreadPool} con un numero fisso di \textit{Worker Thread} che incrementano un contatore condiviso in maniera mutualmente esclusiva.
\begin{lstlisting}[language=Java, numbers=left, stepnumber=1]
import java.util.concurrent.ExecutorService;
import java.util.concurrent.Executors;

public class ThreadPool {
    // Definisco a priori il numero di Task da sottomettere alla Thread Pool
    private static final int NUMBER_OF_TASKS = 10;
    
    public static void main(String[] args) {
        // Inizializzo l'oggetto contatore
        ThreadPoolCounter tpc = new ThreadPoolCounter();
        // Inizializzo la Thread Pool
        ExecutorService executorService = Executors.newFixedThreadPool(5);
        for(int i=0; i<NUMBER_OF_TASKS; i++) {
            // Sottometto alla Thread Pool i Task da eseguire
            executorService.execute(createTask(tpc));
        }
        // Chiudo la Thread Pool
        executorService.shutdown();
    }
    // Creatore di Task
    protected static Runnable createTask(ThreadPoolCounter tpc) {
        // Creo e ritorno un oggetto di Tipo Runnable che esegua l'operazione desiderata
        Runnable runnable = new Runnable() {
            @Override
            public void run() {
                tpc.increaseAndPrintCounter(Thread.currentThread().getName());
            }
        };
        return runnable;
    }
}

class ThreadPoolCounter {
    // Inizializzo una variabile contatore
    private int counter = 0;
    // Metodo esposto ai Thread
    protected synchronized void increaseAndPrintCounter(String threadName) {
        // Incremento il contatore
        counter++;
        // Stampo l'operazione e l'identificativo del Thread che l'ha effettuata
        System.out.println("Il contatore vale " + this.counter + " ed e' stato aggiornato dal Thread " + threadName);
    }
}
\end{lstlisting}
\captionof{lstlisting}{Programma implementante una \textit{Thread Pool}.}
\begin{lstlisting}[language=Bash]
$ java ThreadPool
Il contatore vale 1 ed e' stato aggiornato dal Thread pool-1-thread-1
Il contatore vale 2 ed e' stato aggiornato dal Thread pool-1-thread-5
Il contatore vale 3 ed e' stato aggiornato dal Thread pool-1-thread-4
Il contatore vale 4 ed e' stato aggiornato dal Thread pool-1-thread-4
Il contatore vale 5 ed e' stato aggiornato dal Thread pool-1-thread-4
Il contatore vale 6 ed e' stato aggiornato dal Thread pool-1-thread-3
Il contatore vale 7 ed e' stato aggiornato dal Thread pool-1-thread-2
Il contatore vale 8 ed e' stato aggiornato dal Thread pool-1-thread-4
Il contatore vale 9 ed e' stato aggiornato dal Thread pool-1-thread-5
Il contatore vale 10 ed e' stato aggiornato dal Thread pool-1-thread-1
\end{lstlisting}
\captionof{lstlisting}{Esecuzione del Programma.}
\subsection{Considerazioni}
Nonostante l'utilizzo di \textit{Thread Pool} vada di fatto a diminuire il grado di parallelismo tra Processi Interagenti virtualmente ottenibile, il minore \textit{overhead} causato dal riutilizzo ripetuto dei \textit{Worker Thread} lo rende una soluzione molto più performante della controparte, e di conseguenza spesso utilizzata in applicazioni reali.